% Generated draft: requires mathematical and source review.
\documentclass[11pt]{article}
\usepackage[margin=25mm]{geometry}
\usepackage{fontspec}
\setmainfont{Latin Modern Roman}
\usepackage{amsmath,amssymb,mathtools,unicode-math}
\setmathfont{Latin Modern Math}
\usepackage{enumitem,needspace,xurl,hyperref}
\setlength{\parindent}{0pt}
\setlength{\parskip}{5pt}
\setlength{\emergencystretch}{3em}
\newcommand{\N}{\mathbb{N}}
\newcommand{\Z}{\mathbb{Z}}
\newcommand{\Q}{\mathbb{Q}}
\newcommand{\R}{\mathbb{R}}
\newcommand{\C}{\mathbb{C}}
\newcommand{\F}{\mathbb{F}}
\newcommand{\A}{\mathbb{A}}
\newcommand{\PP}{\mathbb{P}}
\newcommand{\E}{\mathbb{E}}
\newcommand{\eps}{\varepsilon}
\newcommand{\dd}{\,\mathrm{d}}
\begin{document}
\section*{UnsolvedMath ID: 20002833}
\textbf{Source problem code:} catalogue entry
\subsection{Definitions and mathematical statement}
\textbf{Conjecture 3 (Emerton), with source repairs marked.} Let \(\bar\rho:G_\Q\to\operatorname{GL}_2(k)\) be continuous, absolutely irreducible and odd, over a splitting field \(k\) of characteristic \(p\). In the usual \(p>2\) convention, odd means determinant \(-1\) at complex conjugation. Let the level Hecke ideal \(\mathfrak m_N\) have the good-prime eigenvalues \(T_\ell=\operatorname{tr}\bar\rho(\operatorname{Frob}_\ell)\). The precise excluded primes and any determinant/diamond conditions must be supplied. Restoring dropped overlines and subscripts, the conjectural display is of the form
\[
 \varinjlim_N H^1_{\mathrm{\acute et}}(X(N)_{\overline\Q},k)[\mathfrak m_N]
 \ \cong\ \left(\bigotimes_\ell'\pi_\ell^{\mathrm{modified}}
 (\bar\rho|_{G_{\Q_\ell}})\right)\otimes_k\bar\rho.
\]
The local factors are said to be finite-length admissible smooth \(\operatorname{GL}_2(\Q_\ell)\)-representations. Their actual modified definitions, restricted-tensor reference vectors and full action/level conventions are absent. Eigenspace annihilation by \(\mathfrak m\), localization and generalized eigenspaces are different operations; the bracket notation is not silently changed among them.

\textbf{Reported theorem qualification.} The aid attributes the Galois multiplicity-space form to Emerton (2011), Theorem 1.2.6, for \(p>2\) with two explicit exclusions on \(\bar\rho|_{G_{\Q_p}}\). Those exclusions are not stated in this supplied record, so the theorem is cited only conditionally on its full hypotheses, not for all \(p>2\) representations. A multiplicity-space theorem for \(\operatorname{Hom}_{G_\Q}(\bar\rho,-)\) is not automatically the full displayed tensor decomposition in the presence of Galois extensions. The aid mentions later improvements but does not establish unrestricted \(p=2\) or reducible-global versions from its checked sources.

\textbf{Reported finite-level product reduction.} Assume the proven factorization applies to a multiplicity space \(M\), together with the required comparison of finite-level cohomology with invariants at a factorizable neat level \(K=\prod_\ell K_\ell\). Assume also that, almost everywhere, the restricted-tensor reference line is the full one-dimensional \(K_\ell^0\)-fixed space. Then the aid's finite-support tensor-invariants lemma gives
\[
 M^K\cong\bigotimes_\ell'\pi_\ell^{K_\ell},\qquad
 \mu(K):=\dim_kM^K=\prod_\ell d_\ell(K_\ell),
 \quad d_\ell(K_\ell)=\dim_k\pi_\ell^{K_\ell}.
\]
Admissibility gives finite local dimensions and almost all factors are one. The finite-level comparison is a separate hypothesis, especially in characteristic \(p\); neatness or a direct-limit factorization alone does not justify exactness of fixed vectors. If a finite-level eigenspace is additionally known to equal \(M^K\otimes\bar\rho\), its dimension is \(2\mu(K)\); that extra identification is not built into the multiplicity definition.

For two such levels differing only at a prime \(q\), the reported division-free diagnostic is
\[
 \mu(K')\,d_q(K_q)=\mu(K)\,d_q(K'_q).
\]
It remains valid when a local dimension or a common factor vanishes; a quotient formula would not. No explicit \(p\)-local invariant dimensions are computed from missing local factor data.

\textbf{Preserved appended questions.} Calegari asks to calculate finite-level Hecke-eigenspace multiplicities from the modified local factors; Emerton's workshop comment calls the away-from-\(p\) calculation ``easy'' and the \(p\)-part a thesis question. Diamond asks for reducible representations and analogues for Shimura curves attached to quaternion algebras unramified at \(p\), and whether these clarify the local \(p\)-behavior. The final ``3.2. Berger. In the notation of Breuil--Berger'' is an incomplete next-item heading, not a recoverable new theorem. These extensions are not answered by the conditional product lemma. All statements retain the source notes' transcription-error warning and require independent review.
\subsection{Short English statement}
The aid reports Emerton's multiplicity-space theorem only under its missing local hypotheses, then gives a conditional finite-level product of local fixed-vector dimensions. The cross-multiplied one-prime identity survives zero dimensions. The original reducible, Shimura-curve and explicit \(p\)-local questions remain separate, and eigenspaces are not silently identified with localized multiplicity spaces.
\subsection{Sources}
\begin{itemize}
\item Local UnsolvedMath registry, permanent ID \(20002833\), catalogue entry, \texttt{aim-representation-theory-notes.json}, \texttt{notes[104]}, 2006 \emph{\(p\)-adic representations, modularity, and beyond}, Michael Volpato's notes, Conjecture 3 (Emerton), Calegari 3.1.1 and Diamond 3.1.2: \url{https://aimath.org/WWN/padicmodularity/padicmodularity.pdf}.
\item Supplied research attribution: Emerton (2011), Theorem 1.2.6, with \(p>2\) and two local exclusions; no title, full bibliography or exclusions are supplied. Later local-correspondence improvements are mentioned without exact references.
\item Supplied tensor-invariants reduction: attempt 1, partial\_result, partially\_solved, rigor ``proved'', low novelty confidence. The machine-generated aid is not peer reviewed and requires independent verification before citation. The original workshop notes warn of possible transcription errors; no PDF visit is claimed by this draft.
\end{itemize}
\end{document}
